% =====================================================================
%  PeerJ submission — EEG–Hemodynamic Coupling (Topic 9, null result)
%  Class: wlpeerj.cls (Overleaf/PeerJ official template, v1.1)
%  Ported from NCCI_manuscript_draft_v5.tex on 2026-06-11.
%  Citations are now author-year (apalike, set by the class).
%  Figures live one directory up (manuscript/), hence \graphicspath.
% =====================================================================
\documentclass[fleqn,10pt,lineno]{wlpeerj}

\usepackage{float}      % for [H] placement
\usepackage{array}
\graphicspath{{../}}    % figures are in ../ (manuscript/)

\usepackage[breaklinks=true]{hyperref}
\definecolor{OldCyan}{RGB}{62,127,137}
\hypersetup{colorlinks,citecolor=OldCyan,filecolor=OldCyan,
            linkcolor=OldCyan,urlcolor=OldCyan,
            pdftitle={EEG-Hemodynamic Coupling During General Anaesthesia:
                      A Large-Scale Exploratory Analysis in 2,070 Surgical Cases
                      with Null Acute Event-Response Findings},
            pdfauthor={Ge Gao},
            pdfkeywords={EEG-hemodynamic coupling; Granger causality;
                         cross-wavelet coherence; general anaesthesia; null result; VitalDB}}

% PeerJ does not number sections
\setcounter{secnumdepth}{0}

\title{EEG--Hemodynamic Coupling During General Anaesthesia:
A Large-Scale Exploratory Analysis in 2,070 Surgical Cases
with Null Acute Event-Response Findings}

\author[1]{Ge Gao}

\affil[1]{Department of Anesthesiology, The First Affiliated Hospital,
Zhejiang University School of Medicine, Hangzhou, China}
\affil[ ]{\textit{Corresponding author:} Ge Gao, gaoge2018@zju.edu.cn}

\keywords{EEG--hemodynamic coupling, Granger causality, cross-wavelet coherence,
general anaesthesia, intraoperative monitoring, null result, VitalDB, exploratory analysis}

\begin{abstract}
\textbf{Background.}
Whether dynamic coupling between cortical EEG and arterial blood pressure (ABP) varies
systematically across patients and responds to acute haemodynamic excursions during
general anaesthesia has not been characterised at large scale. We report an exploratory
retrospective analysis with predominantly null findings.

\textbf{Methods.}
We analysed 2,070 general anaesthesia cases from VitalDB with simultaneous raw EEG
(128~Hz) and invasive ABP (500~Hz). Granger causality (GC) and cross-wavelet coherence
(WC) were computed in 10-minute sliding windows. A case-level composite summary score
(the Neuro-Cardiovascular Coupling Index, NCCI) was derived by global PCA. A per-case
Window Coupling Score (WCS), sign-aligned and MinMax-normalised within each patient, was
used for event analysis. Surrogate haemodynamic excursions were defined as MAP rise
$>$20~mmHg within 90~s with BIS~$<$60. The primary statistical analysis was a generalised
estimating equation (GEE) on per-event $\Delta$WCS (exchangeable; sensitivity: AR(1)) to
account for within-case clustering.

\textbf{Results.}
EEG changes temporally preceded MAP changes in 71\% of cases. Mean NCCI was
0.727~$\pm$~0.067. Surrogate excursions were detected in 2,031/2,070 cases (98.1\%),
indicating non-specific endpoint capture. Primary result (GEE-Exchangeable):
$\beta = +0.0038$ (95\% CI $+0.0001$ to $+0.0076$; $p = 0.047$), Cohen's $d = 0.047$
(negligible). With the more conservative AR(1) correlation structure $\beta = +0.0020$,
$p = 0.230$ (not significant). Stricter excursion thresholds yielded null or negative
effects across all definitions.

\textbf{Conclusions.}
There is no robust evidence of WCS change at surrogate haemodynamic excursions. The
predominantly null result reflects three compounding limitations: a non-specific surrogate
endpoint (98.1\% case rate), a mismatch between 10-minute coupling windows and 90-second
excursions, and pharmacodynamic confounding uncontrolled in the current data. These
findings are hypothesis-generating and highlight the methodological requirements for future
coupling research with proper event labels and drug-infusion covariates.
\end{abstract}

\begin{document}

\flushbottom
\maketitle
\thispagestyle{empty}

% ── INTRODUCTION ─────────────────────────────────────────────────────────────
\section{Introduction}

During general anaesthesia, the brain and cardiovascular system undergo coupled physiological
changes driven by anaesthetic agents \citep{brown2010,purdon2015}. EEG monitoring via the
Bispectral Index (BIS) is widely used to track anaesthetic depth \citep{avidan2008}, while
arterial blood pressure (ABP) provides continuous haemodynamic surveillance. Whether changes
in EEG activity temporally precede or follow haemodynamic changes --- and whether their
coupling strength varies meaningfully across patients --- are clinically relevant empirical
questions that have received limited systematic attention at scale.

Granger causality (GC) \citep{granger1969} is a widely applied method for testing temporal
precedence between physiological time series \citep{pereda2005,seth2015}. Cross-wavelet
coherence (WC) \citep{torrence1998} provides frequency-resolved coupling without requiring
stationarity, and has been used in cardiovascular--neurological coupling studies
\citep{pikovsky2001}. The Alpha/Delta Ratio (ADR), derived from frontal EEG power spectra,
provides a continuous index of cortical arousal that is sensitive to anaesthetic depth
\citep{akeju2014,bruhn2000}.

The VitalDB open dataset \citep{lee2022} provides 6,388 surgical cases with synchronised EEG
and ABP waveforms, offering a unique opportunity to characterise EEG--hemodynamic coupling
systematically. Prior studies of intraoperative coupling have been limited to small samples
or lacked the multivariate signal processing infrastructure available here.

Motivated by this gap, we conducted a planned exploratory retrospective analysis with
three objectives: (1) to quantify the direction and magnitude of EEG--MAP temporal precedence
(GC) and spectral coherence (WC) across a large surgical cohort; (2) to construct a case-level
composite coupling summary score (NCCI) and characterise its population distribution; and
(3) to test whether a per-case window-level coupling score (WCS) is associated with surrogate
haemodynamic excursions, defined by MAP and BIS criteria. \textbf{This study makes no clinical
validation claims and is explicitly hypothesis-generating.} The study is reported in accordance
with the STROBE guideline for observational studies; the completed checklist is provided as
Supplementary material.

% ── MATERIALS & METHODS ──────────────────────────────────────────────────────
\section{Materials \& Methods}

\subsection{Study Design and Data Source}

Retrospective observational cohort study using VitalDB v1.0.0 \citep{lee2022}, a
high-resolution perioperative dataset of patients who underwent surgery at Seoul National
University Hospital (Seoul, Republic of Korea) between August 2016 and June 2017. The original
data collection was approved by the Seoul National University Hospital Institutional Review
Board (H-1408-101-605), which waived the requirement for informed consent owing to the use of
fully anonymised records. The dataset is openly available on PhysioNet under the Creative
Commons Attribution 4.0 International (CC BY 4.0) licence (version DOI:
10.13026/czw8-9p62); the present secondary analysis of anonymised public data required no
additional ethical approval.

\subsection{Patient Selection}

Inclusion criteria, applied to all 6,388 cases in VitalDB v1.0.0, were: age $\geq$18 years;
general anaesthesia; operation duration $\geq$60~min; and the presence of both raw EEG
(128~Hz) and invasive ABP (500~Hz) waveforms with $\geq$30~min of concurrent recording. Cases
from cardiac or thoracic surgical departments were excluded; no additional diagnosis-based
clinical exclusions were applied. Signal-level artefact handling (EEG amplitude $>$200~$\mu$V
and out-of-range ABP) is described in the next section. The selection flowchart
(Figure~\ref{fig:flowchart}) proceeds from 6,388 cases $\to$ 3,825 after the basic clinical
filters (age~$\geq$18, general anaesthesia, non-cardiac/non-thoracic department, operation
$\geq$60~min) $\to$ 2,076 with both EEG and ABP waveforms present $\to$ 2,070 after requiring
$\geq$30~min of concurrent valid EEG and ABP signal.

\subsection{Signal Processing}

\textbf{EEG} (128~Hz): Bandpass 0.5--45~Hz (4th-order Butterworth) + 50~Hz notch. Segments
$>$200~$\mu$V masked. Per 30-second window (10-second step): Welch PSD $\to$
$\delta/\theta/\alpha/\beta$ band powers, Alpha/Delta Ratio (ADR), SEF95, spectral entropy,
suppression ratio.

\textbf{ABP} (500~Hz): Beat-to-beat peak detection (min height 20~mmHg; min inter-peak 200
samples). Per window: MAP, SBP/DBP, pulse pressure, PPV, SDNN, RMSSD, dP/dt$_\text{max}$.

ADR was selected as the primary EEG coupling variable on three grounds: (i)~prior evidence
that ADR tracks BIS-defined anaesthetic depth \citep{bruhn2000,akeju2014}; (ii)~it is a
continuous, interpretable scalar at each window; (iii)~it showed the lowest artefact-driven
outlier rate among tested features. MAP was selected as the primary haemodynamic variable
because it integrates systolic and diastolic components and is the principal clinical target
for haemodynamic management.

\subsection{Granger Causality}

Pairwise GC between ADR and MAP in 60-window (10-minute) sliding windows (3-window step)
\citep{granger1969,barnett2014}. Stationarity: ADF test ($\alpha = 0.05$); non-stationary
series first-differenced. Lag: 1--5, selected by minimum SSR F-test $p$-value. GC$_\text{e2h}$
was right-skewed (raw mean 2.61, SD 5.08) and winsorised at the 99th percentile
($\hat{c} = 4.26$) for PCA input; raw values were used for directionality. GC quantifies
temporal precedence, not mechanistic causation \citep{granger1969}.

\subsection{Cross-Wavelet Coherence}

Continuous wavelet transform with a complex Morlet wavelet (cmor1.5-1.0, $dt = 10$~s)
\citep{torrence1998}. Coherence:
\begin{equation}
C_{xy}(a,t) = \frac{|\langle W_x W_y^* \rangle|^2}
  {\langle |W_x|^2 \rangle \cdot \langle |W_y|^2 \rangle}
\end{equation}
Smoothed with a uniform filter ([5,~11]). Cone-of-influence (COI) edge effects affect
$\approx$5\% of data at record boundaries; a supplementary sensitivity analysis
(Supplementary S5) confirmed that removing COI-affected samples did not materially alter
WC$_\text{overall}$ (mean change $<$0.003). Summary metrics: WC$_\text{overall}$,
WC$_\text{low}$ ($<$0.01~Hz), WC$_\text{mid}$ (0.01--0.03~Hz).

\subsection{Coupling Scores: NCCI and WCS}

\textbf{Case-level NCCI (Neuro-Cardiovascular Coupling Index)} (one value per case): Global PCA across all 2,070 cases on per-case
means of [GC$_\text{e2h}$ F (winsorised), WC$_\text{overall}$, Pearson $r$]. $Z$-score
standardised; PC1 MinMax to [0,1]. Sign convention: PC1 inverted if corr(PC1,~WC)$<$0
(ensuring NCCI~$=$~1 $\equiv$ maximum WC and GC).

\textbf{Window Coupling Score (WCS)} (per-window; for event analysis only): Per-case local PCA
on per-window [GC$_\text{e2h}$ F, WC, Pearson $r$]. The same sign convention is applied per
case. PC1 scores are MinMax-normalised within each case. WCS~$=$~1 consistently represents
higher WC and GC within that patient's own range. WCS values cannot be compared across
patients and are not on the same scale as NCCI.

\subsection{Surrogate Haemodynamic Excursions}

MAP rise $>$20~mmHg within 90~s while BIS~$<$60. Pre-event window: the 10-minute window
immediately preceding the excursion. Post-event window: the 10-minute window immediately
following. This definition has known limitations: no surgical timestamps; MAP rises may
reflect non-nociceptive causes; the threshold is empirically chosen. We therefore term these
\textit{surrogate haemodynamic excursions}, not nociceptive events.

\subsection{Statistical Analysis}

\textbf{Primary --- GEE on $\Delta$WCS (Exchangeable).} For each excursion, $\Delta$WCS
$=$ WCS$_\text{post}$ $-$ WCS$_\text{pre}$. A GEE with exchangeable within-case correlation
\citep{liang1986} (formula: $\Delta$WCS $\sim$ 1, the intercept tests departure from zero),
robust sandwich variance, grouped by patient ($n = 2{,}031$ patients, 73,001 observations).
This directly models the within-case relative change without requiring WCS values to be
comparable across patients.

\textbf{Sensitivity --- GEE-AR(1).} The same model with an AR(1) working correlation
structure, which accounts for the temporal ordering of events within each case (more
conservative for time-series data).

\textbf{Secondary --- patient-level.} Per-patient mean $\Delta$WCS tested against zero
(one-sample $t$-test on patient means, $n = 2{,}031$). Effect size: Cohen's $d$.

\textbf{Threshold sensitivity.} The primary GEE analysis was repeated for four stricter
excursion definitions (Table~\ref{tab:threshold}).

Continuous variables are reported as mean~$\pm$~SD or median [IQR]. Two-sided $\alpha = 0.05$.
Software: Python 3.10 (NumPy 1.24, SciPy 1.11, statsmodels 0.14, PyWavelets 1.4,
scikit-learn 1.3).

% ── RESULTS ──────────────────────────────────────────────────────────────────
\section{Results}

\subsection{Cohort}

2,070 cases were analysed (Figure~\ref{fig:flowchart}; Table~\ref{tab:cohort}): median
anaesthesia duration 262~[195--340]~min; age 58.3~$\pm$~14.5 years; 57\% male; ASA I--II in
82\%. A representative single-case recording, showing the synchronised EEG band powers, BIS,
haemodynamics, and cross-wavelet coherence used throughout the analysis, is shown in
Figure~\ref{fig:casepanel}.

\subsection{Granger Causality}

Mean GC$_\text{e2h}$ F (winsorised) was 2.49~$\pm$~0.50; GC$_\text{h2e}$ was 2.23~$\pm$~0.50
(Supplementary Figure~S1).
EEG temporally preceded MAP in 71\% of cases (GC directionality, winsorised:
0.28~$\pm$~0.68). \textit{Interpretation caveat:} given the high low-frequency WC (below),
shared anaesthetic pharmacodynamic drift is a plausible primary driver of this temporal
pattern; no mechanistic claim is made.

\subsection{Cross-Wavelet Coherence}

WC$_\text{overall}$ was 0.799~$\pm$~0.034. WC$_\text{low}$ ($<$0.01~Hz) was 0.824~$\pm$~0.033,
reflecting shared slow pharmacodynamic trends (periods $>$100~s) most likely from propofol and
remifentanil co-variation. WC$_\text{mid}$ (0.01--0.03~Hz) was 0.352~$\pm$~0.025.
Representative heatmaps for three NCCI tertile cases are shown in Figure~\ref{fig:wavelet}.

\subsection{Case-Level NCCI Distribution}

NCCI was 0.727~$\pm$~0.067 (median 0.729; IQR 0.694--0.766; 5th--95th pct [0.624, 0.827]).
PC1 loadings: GC 0.65, WC 0.73, Pearson $-$0.23; PC1 variance 37.0\%, indicating limited
shared information among the three input metrics. The distribution is shown in
Figure~\ref{fig:distributions}. Quartile stratification (Supplementary Figure~S2)
shows a monotonic component ordering that is a \textit{construction artefact} of the NCCI
definition and has no independent evidentiary value.

\subsection{WCS Response to Surrogate Haemodynamic Excursions}

Surrogate excursions were detected in 2,031/2,070 cases (98.1\%), yielding 73,001 observations
(median 29 per case [IQR 18--45]; max 212). The 98.1\% case rate indicates this definition is
non-specific and cannot discriminate nociceptive from non-nociceptive haemodynamic variation.

\textbf{Primary result (GEE-Exchangeable, $n = 2{,}031$ patients):} $\beta = +0.0038$
(SE $= 0.0019$, 95\% CI $+0.0001$ to $+0.0076$, $z = 1.99$, $p = 0.047$). Cohen's
$d = 0.047$ (negligible by conventional standards).

\textbf{GEE-AR(1) sensitivity:} $\beta = +0.0020$ (SE $= 0.0017$, 95\% CI $-0.0013$ to
$+0.0054$, $p = 0.230$). The shift from $p = 0.047$ (exchangeable) to $p = 0.230$ (AR(1))
when the temporal ordering of events is accounted for indicates that the borderline result in
the primary analysis is \textbf{not robust to the correlation structure choice}.

\textbf{Patient-level secondary analysis:} Mean $\Delta$WCS per patient was
$+0.0047~\pm~0.0997$; pre-excursion 0.333~$\pm$~0.124; post-excursion 0.337~$\pm$~0.131.
The difference is statistically significant but of negligible magnitude and inconsistent with
the primary GEE-AR(1) estimate, likely reflecting high precision from a large sample size
rather than a meaningful effect ($d = 0.047$). 49\% of patients showed a negative mean
$\Delta$WCS.

\textbf{Cluster bootstrap:} $\Delta$WCS $= +0.0047$; 95\% CI [$+0.0004$, $+0.0090$]. The
lower bound is near zero, consistent with GEE-AR(1).

\textbf{Threshold sensitivity} (Table~\ref{tab:threshold}): All four stricter definitions
yield case rates $>$95\%. The MAP$\geq$25 threshold yields $\beta \approx 0$ ($p = 0.976$);
MAP$\geq$30 thresholds yield negative $\beta$ values. No threshold produced a clinically
meaningful effect ($|d| \leq 0.05$). Figure~\ref{fig:excursions} summarises the event analysis;
the case-level Alpha/Delta Ratio analogue is shown in Supplementary Figure~S3.

% ── DISCUSSION ───────────────────────────────────────────────────────────────
\section{Discussion}

\subsection{Principal Findings}

This large exploratory study of EEG--hemodynamic coupling in 2,070 general anaesthesia cases
yields one positive descriptive finding and one primary negative finding. Descriptively, EEG
changes temporally precede MAP changes by GC in 71\% of cases, and overall coupling coherence
is high (WC $= 0.799$), predominantly at low frequencies consistent with shared pharmacodynamic
co-variation. The negative finding: WCS shows no robust change at surrogate haemodynamic
excursions. The GEE-Exchangeable result ($p = 0.047$) is borderline but shifts to
non-significant ($p = 0.230$) with AR(1) correlation --- a more appropriate structure for
time-ordered events within patients. Effect sizes are negligible ($d = 0.047$) across all
analyses.

\subsection{The Endpoint Defines the Finding}

We propose that \textbf{the primary negative result is more informative about the inadequacy of
the surrogate excursion definition than about the absence of true EEG--hemodynamic coupling
dynamics.} A 98.1\% excursion rate means the definition captures nearly universal haemodynamic
variation --- fluid administration, positional change, autonomic fluctuations --- alongside any
genuine nociceptive responses. In such a setting, any acute coupling metric will appear
insensitive, not because coupling dynamics are absent, but because the endpoint cannot isolate
the physiological event of interest. The threshold sensitivity analysis reinforces this:
stricter MAP/BIS criteria continue to capture $>$95\% of cases, indicating the problem is
structural rather than threshold-dependent.

\subsection{Window--Event Mismatch}

A 10-minute coupling window is well-suited to characterising slow coupling trends but is
fundamentally mismatched to 90-second haemodynamic excursions. The WCS likely captures propofol
and remifentanil pharmacodynamic co-variation at the 100--600-second timescale
\citep{struys2000} rather than rapid autonomic responses. Future work should explore coupling
metrics derived from shorter windows ($<$2 minutes) or use event-locked analysis approaches
rather than sliding windows.

\subsection{GC Temporal Precedence}

The 71\% EEG-preceding result is descriptively noteworthy but difficult to interpret causally.
Without drug infusion covariates, the temporal structure of propofol and remifentanil
concentration trajectories may introduce apparent EEG-to-MAP precedence through shared
pharmacokinetics. Whether a true neural coupling pathway (e.g.\ cortical suppression $\to$
reduced sympathetic output $\to$ MAP reduction) contributes cannot be disentangled from the
available data \citep{sessler2012}.

\subsection{Value of Negative Results in Physiological Signal Research}

Negative and null results from large open datasets constitute scientifically valuable
contributions, even when they do not yield immediately actionable clinical findings. This study
demonstrates: (i)~that EEG--MAP coupling can be systematically quantified at population scale
using VitalDB; (ii)~that the coupling composite score NCCI captures population-level variation;
and (iii)~that the MAP/BIS-based surrogate excursion definition is insufficiently specific for
acute coupling validation in a retrospective setting. This third finding --- that the
\textit{endpoint} fails before the \textit{metric} can be properly tested --- is a
methodological contribution to future intraoperative monitoring research
\citep{sessler2012,struys2000}. Researchers designing future coupling studies should use
surgical stimulus timestamps, validated pain scales, or drug-intervention event logs as
endpoints, rather than haemodynamic thresholds alone.

\subsection{Limitations}

(1) Single-centre; generalisability unknown. (2) Frontal single-channel EEG only. (3) No
surgical stimulus timestamps; the 98.1\% surrogate excursion rate indicates a non-specific
endpoint. (4) Drug infusion rates not included as covariates; pharmacodynamic co-variation
likely dominates the measured coupling. (5) GC assumes linearity; the lag was selected by
minimum $p$-value (which may bias toward significance; fixed-lag sensitivity analyses are
reported in Supplementary S3 and recommended for future work). (6) NCCI and WCS are distinct
constructs; WCS values are not comparable across patients. (7) 10-minute WCS windows cannot
track 90-second excursions. (8) No external validation cohort.

% ── CONCLUSIONS ──────────────────────────────────────────────────────────────
\section{Conclusions}

In 2,070 general anaesthesia cases, EEG temporally preceded MAP changes in 71\% of cases and
overall coherence was high, predominantly reflecting shared pharmacodynamic slow trends. A
case-level composite coupling score (NCCI) captured systematic inter-patient variation but was
self-referentially constructed.

The \textbf{primary finding is null:} WCS shows no robust change at surrogate haemodynamic
excursions (GEE-Exchangeable $p = 0.047$; GEE-AR(1) $p = 0.230$; $d = 0.047$). This result
reflects three compounding limitations: endpoint non-specificity (98.1\% case rate),
window--event timescale mismatch, and uncontrolled pharmacodynamic confounding.

The null finding has methodological value: it demonstrates that MAP/BIS-based surrogate
excursion definitions are insufficient for acute coupling validation in retrospective
intraoperative data. Future coupling studies should incorporate surgical timestamps, shorter
coupling windows matched to event timescales, and drug infusion covariates.

% =====================================================================
%  ADDITIONAL INFORMATION AND DECLARATIONS  (PeerJ-required block)
% =====================================================================
\section{Additional Information and Declarations}

\subsection{Ethics}
This study is a secondary analysis of the VitalDB v1.0.0 open dataset. The original data
collection was approved by the Institutional Review Board of Seoul National University Hospital
(H-1408-101-605), which waived the requirement for informed consent owing to the use of fully
anonymised records. All records are publicly available on PhysioNet under the Creative Commons
Attribution 4.0 International (CC BY 4.0) licence; no additional ethical approval was required
for the present analysis.

\subsection{Funding}
The author received no specific funding for this work.

\subsection{Competing Interests}
The author declares that there are no competing interests.

\subsection{Author Contributions}
\textbf{Ge Gao} conceived and designed the study, performed the experiments, analysed the data,
prepared the software, and authored and approved the final manuscript.

\subsection{Data Availability}
VitalDB v1.0.0 is openly available on PhysioNet
(\url{https://physionet.org/content/vitaldb/1.0.0/}; version DOI: 10.13026/czw8-9p62) under the
Creative Commons Attribution 4.0 International (CC BY 4.0) licence. No additional data were
collected.

\subsection{Code Availability}
The full analysis pipeline (Python 3.10; NumPy, SciPy, statsmodels, PyWavelets, scikit-learn)
is openly available at \url{https://github.com/ZJUArtemis/eeg-hemodynamic-coupling}. Key
scripts include signal processing, GC and WC computation, NCCI/WCS construction, and
statistical analysis.

\subsection{Supplemental Information}
Supplementary text S1--S6 and Supplementary Figures S1--S3 are provided as separate files.

% ── REFERENCES ──────────────────────────────────────────────────────────────
\bibliography{references}

% =====================================================================
%  TABLES
% =====================================================================
\clearpage

\begin{table}[H]
\centering
\caption{Cohort characteristics, coupling metrics, and WCS event analysis ($n = 2{,}070$).}
\label{tab:cohort}
\begin{tabular}{lc}
\toprule
\textbf{Variable} & \textbf{Value} \\
\midrule
\multicolumn{2}{l}{\textit{Demographics and clinical}} \\
Age (mean $\pm$ SD, years)        & 58.3 $\pm$ 14.5 \\
Male sex                          & 1,180 (57.0\%) \\
BMI (mean $\pm$ SD, kg/m$^2$)     & 23.0 $\pm$ 3.7 \\
ASA I / II / III / IV             & 458 / 1,201 / 342 / 24 \\
Anaesthesia duration (median [IQR], min) & 262 [195--340] \\
\midrule
\multicolumn{2}{l}{\textit{Granger causality (mean $\pm$ SD)}} \\
GC$_\text{e2h}$ F (winsorised)    & 2.49 $\pm$ 0.50 \\
GC$_\text{e2h}$ F (raw)           & 2.61 $\pm$ 5.08 \\
GC$_\text{h2e}$ F                 & 2.23 $\pm$ 0.50 \\
EEG preceding MAP                 & 71\% of cases \\
Pearson $r$ (ADR vs MAP)          & 0.036 $\pm$ 0.181 \\
\midrule
\multicolumn{2}{l}{\textit{Cross-wavelet coherence (mean $\pm$ SD)}} \\
WC$_\text{overall}$               & 0.799 $\pm$ 0.034 \\
WC$_\text{low}$ ($<$0.01 Hz)      & 0.824 $\pm$ 0.033 \\
WC$_\text{mid}$ (0.01--0.03 Hz)   & 0.352 $\pm$ 0.025 \\
\midrule
\multicolumn{2}{l}{\textit{NCCI (case-level composite summary score)}} \\
Mean $\pm$ SD                     & 0.727 $\pm$ 0.067 \\
Median [IQR]                      & 0.729 [0.694--0.766] \\
5th--95th pct                     & [0.624, 0.827] \\
PCA PC1 variance                  & 37.0\% \\
PC1 loadings                      & GC 0.65, WC 0.73, Pearson $-$0.23 \\
\midrule
\multicolumn{2}{l}{\textit{WCS excursion analysis (per-case local PCA)$^{\dagger}$}} \\
Cases with $\geq$1 excursion      & 2,031/2,070 (98.1\%) \\
Total excursions (observations)   & 73,001 \\
Excursions/case: median [IQR]     & 29 [18--45] \\
Pre-excursion WCS (patient mean $\pm$ SD)  & 0.333 $\pm$ 0.124 \\
Post-excursion WCS (patient mean $\pm$ SD) & 0.337 $\pm$ 0.131 \\
Patient-mean $\Delta$WCS $\pm$ SD & $+0.0047 \pm 0.0997$ \\
GEE-Exchangeable $\beta$ (95\% CI; $p$) & $+0.0038$ ($+0.0001$ to $+0.0076$; $p = 0.047$) \\
GEE-AR(1) $\beta$ (95\% CI; $p$)  & $+0.0020$ ($-0.0013$ to $+0.0054$; $p = 0.230$) \\
Cohen's $d$                       & 0.047 (negligible) \\
\% patients negative mean $\Delta$WCS & 49\% \\
\bottomrule
\end{tabular}

\vspace{4pt}
\footnotesize $^{\dagger}$WCS (per-case local PCA, sign-aligned, MinMax within case) is a
distinct construct from NCCI (global PCA). WCS values are not comparable across patients. GEE
analyses model $\Delta$WCS per excursion clustered by patient.
\end{table}

\begin{table}[H]
\centering
\caption{Surrogate excursion definition sensitivity analysis.}
\label{tab:threshold}
\small
\begin{tabular}{p{4.2cm}rrp{3.0cm}rr}
\toprule
\textbf{Definition} & \textbf{Cases (\%)} & \textbf{Events} & \textbf{Mean $\Delta$WCS (95\% CI)} & \textbf{Cohen's $d$} & \textbf{$p$} \\
\midrule
\textbf{MAP$\geq$20, BIS$<$60} & \textbf{2,031 (98.1\%)} & \textbf{73,001} & \textbf{$+0.0038$ [$+0.0001$, $+0.0076$]} & \textbf{+0.047} & \textbf{0.047} \\
MAP$\geq$25, BIS$<$60 & 2,013 (97.2\%) & 55,699 & $-0.0001$ [$-0.0043$, $+0.0041$] & +0.008 & 0.976 \\
MAP$\geq$30, BIS$<$60 & 1,990 (96.1\%) & 45,502 & $-0.0040$ [$-0.0086$, $+0.0006$] & $-$0.026 & 0.089 \\
MAP$\geq$20, BIS$<$55 & 2,026 (97.9\%) & 69,616 & $+0.0033$ [$-0.0005$, $+0.0072$] & +0.042 & 0.085 \\
MAP$\geq$30, BIS$<$55 & 1,979 (95.6\%) & 43,539 & $-0.0043$ [$-0.0089$, $+0.0004$] & $-$0.026 & 0.072 \\
\bottomrule
\end{tabular}

\vspace{4pt}
\footnotesize GEE-Exchangeable on $\Delta$WCS, clustered by patient. None of the stricter
definitions reaches a clinically meaningful effect magnitude ($|d| \leq 0.05$). MAP$\geq$25 and
MAP$\geq$30 thresholds produce near-zero or negative $\beta$ values. Case rates $>$95\% across
all definitions confirm the non-specificity of the surrogate endpoint.
\end{table}

% =====================================================================
%  FIGURES
% =====================================================================
\clearpage

\begin{figure}[H]
\centering
\includegraphics[width=0.72\linewidth]{fig1_flowchart.png}
\caption{Patient selection flowchart. Exclusion steps from 6,388 VitalDB cases to the 2,070
cases in the coupling analysis and 2,031 cases contributing to the WCS excursion analysis.}
\label{fig:flowchart}
\end{figure}

\begin{figure}[H]
\centering
\includegraphics[width=\linewidth]{fig2_case_panel.png}
\caption{Representative intraoperative time series for a single case. \textbf{(A)}~EEG delta and
alpha band powers. \textbf{(B)}~BIS value. \textbf{(C)}~MAP and HR. \textbf{(D)}~Cross-wavelet
coherence heatmap (frequency $\times$ time). \textbf{(E)}~Alpha/Delta Ratio (ADR). Time in
minutes from anaesthesia start.}
\label{fig:casepanel}
\end{figure}

\begin{figure}[H]
\centering
\includegraphics[width=\linewidth]{fig3_wavelet_coherence.png}
\caption{Wavelet coherence heatmaps for three representative cases, from the lowest (top;
NCCI $= 0.674$), middle (centre; NCCI $= 0.730$), and highest (bottom; NCCI $= 0.783$) NCCI
tertile. Colour: coherence 0--1. High low-frequency coherence ($<$0.01~Hz) is visible in all
cases, consistent with shared pharmacodynamic slow drift.}
\label{fig:wavelet}
\end{figure}

\begin{figure}[H]
\centering
\includegraphics[width=\linewidth]{fig4_distributions.png}
\caption{Distribution of case-level coupling metrics ($n = 2{,}070$ cases). \textbf{(A)}~NCCI
(composite summary score, global PCA). \textbf{(B)}~WC$_\text{overall}$. \textbf{(C)}~GC
directionality (clipped to $\pm$10 for display; positive $=$ EEG preceding MAP). Vertical
lines: mean $\pm$ 1~SD.}
\label{fig:distributions}
\end{figure}

\begin{figure}[H]
\centering
\includegraphics[width=\linewidth]{fig5_wcs_excursions.png}
\caption{Window Coupling Score (WCS) at surrogate haemodynamic excursions ($n = 2{,}031$
patients, 73,001 events). \textbf{(A)}~Paired scatter: per-patient mean pre- vs
post-excursion WCS; identity line in grey. \textbf{(B)}~Distribution of per-patient mean
$\Delta$WCS; key statistical results annotated. \textbf{(C)}~MAP rise magnitude vs mean
$\Delta$WCS (decile means; event-level, descriptive only). WCS (per-case local PCA) is not on
the same scale as NCCI (global PCA).}
\label{fig:excursions}
\end{figure}

\end{document}
